\documentclass[12pt,a4paper]{article}

\usepackage[utf8]{inputenc}
\usepackage[T1]{fontenc}
\usepackage[brazil]{babel}
\usepackage{amsmath,amssymb}
\usepackage{geometry}
\usepackage{graphicx}
\usepackage{float}
\usepackage{booktabs}
\usepackage{hyperref}
\usepackage{indentfirst}
\usepackage{listings}
\usepackage{xcolor}

\geometry{margin=2.5cm}

\lstdefinestyle{pythonstyle}{
    language=Python,
    basicstyle=\ttfamily\small,
    keywordstyle=\color{blue},
    commentstyle=\color{gray},
    stringstyle=\color{teal},
    numbers=left,
    numberstyle=\tiny,
    breaklines=true,
    frame=single,
    tabsize=4
}

\title{Trabalho de Camada Limite -- 1/2026\\
Camada limite laminar sobre placa plana}
\author{Gustavo Bittar Gonçalves}
\date{Entrega: 10/07/2026}

\begin{document}

\maketitle

\section{Introdução}

Este relatório apresenta a modelagem física e matemática da camada limite laminar formada sobre uma placa plana, considerando fluido newtoniano, escoamento permanente, incompressível e velocidade uniforme $U$ na entrada. A partir das equações diferenciais parciais de Navier--Stokes, é feita uma análise de ordem de grandeza para obter as equações simplificadas de camada limite. Em seguida, utiliza-se a transformação de similaridade de Blasius para reduzir o problema a uma equação diferencial ordinária, que é resolvida numericamente pelo método de Runge--Kutta de quarta ordem combinado com o método de shooting.

\section{Modelo físico}

Considera-se uma placa plana semi-infinita alinhada com o escoamento externo uniforme. O fluido escoa na direção $x$ com velocidade $U$ longe da parede, enquanto $y$ representa a direção normal à placa. A placa é fixa, lisa e impermeável. O fluido é newtoniano, incompressível e possui viscosidade cinemática constante $\nu$.

As hipóteses físicas adotadas são:

\begin{itemize}
    \item escoamento bidimensional;
    \item regime permanente;
    \item fluido incompressível;
    \item propriedades constantes;
    \item ausência de gradiente de pressão externo, isto é, $\frac{dp}{dx}=0$;
    \item condição de não deslizamento na parede;
    \item velocidade externa uniforme $U$.
\end{itemize}

As condições de contorno físicas são:

\begin{align}
    u(x,0) &= 0, & v(x,0) &= 0, \\
    u(x,y\rightarrow\infty) &= U.
\end{align}

\section{Modelo matemático diferencial simplificado}

Para um escoamento incompressível, bidimensional e permanente, as equações de conservação são dadas por:

\begin{equation}
    \frac{\partial u}{\partial x} + \frac{\partial v}{\partial y} = 0,
\end{equation}

\begin{equation}
    u\frac{\partial u}{\partial x} + v\frac{\partial u}{\partial y}
    = -\frac{1}{\rho}\frac{\partial p}{\partial x}
    + \nu\left(\frac{\partial^2 u}{\partial x^2}+\frac{\partial^2 u}{\partial y^2}\right),
\end{equation}

\begin{equation}
    u\frac{\partial v}{\partial x} + v\frac{\partial v}{\partial y}
    = -\frac{1}{\rho}\frac{\partial p}{\partial y}
    + \nu\left(\frac{\partial^2 v}{\partial x^2}+\frac{\partial^2 v}{\partial y^2}\right).
\end{equation}

Na camada limite, seja $L$ uma escala longitudinal e $\delta$ a espessura da camada limite. Para $\delta \ll L$, tem-se:

\begin{equation}
    u \sim U, \qquad x \sim L, \qquad y \sim \delta.
\end{equation}

Da continuidade:

\begin{equation}
    \frac{U}{L} \sim \frac{V}{\delta}
    \quad \Rightarrow \quad
    V \sim U\frac{\delta}{L}.
\end{equation}

Na equação de quantidade de movimento em $x$:

\begin{align}
    u\frac{\partial u}{\partial x} &\sim \frac{U^2}{L}, \\
    v\frac{\partial u}{\partial y} &\sim \left(U\frac{\delta}{L}\right)\frac{U}{\delta}=\frac{U^2}{L}, \\
    \nu\frac{\partial^2 u}{\partial x^2} &\sim \nu\frac{U}{L^2}, \\
    \nu\frac{\partial^2 u}{\partial y^2} &\sim \nu\frac{U}{\delta^2}.
\end{align}

Como $\delta \ll L$, o termo difusivo na direção $x$ é pequeno em relação ao termo difusivo na direção $y$:

\begin{equation}
    \nu\frac{U}{L^2} \ll \nu\frac{U}{\delta^2}.
\end{equation}

Equilibrando inércia e difusão viscosa normal:

\begin{equation}
    \frac{U^2}{L} \sim \nu\frac{U}{\delta^2}
    \quad \Rightarrow \quad
    \frac{\delta}{L} \sim \frac{1}{\sqrt{Re_L}},
\end{equation}

onde:

\begin{equation}
    Re_L = \frac{UL}{\nu}.
\end{equation}

Assim, para $Re_L$ elevado, a camada limite é fina. O modelo diferencial simplificado fica:

\begin{equation}
    \boxed{\frac{\partial u}{\partial x} + \frac{\partial v}{\partial y} = 0}
\end{equation}

\begin{equation}
    \boxed{u\frac{\partial u}{\partial x} + v\frac{\partial u}{\partial y}
    = \nu\frac{\partial^2 u}{\partial y^2}}
\end{equation}

com as condições:

\begin{align}
    u(x,0)&=0, & v(x,0)&=0, & u(x,\infty)&=U.
\end{align}

\section{Teoria de similaridade}

Para reduzir as equações diferenciais parciais a uma equação diferencial ordinária, introduz-se a função corrente $\psi$:

\begin{equation}
    u = \frac{\partial \psi}{\partial y}, \qquad
    v = -\frac{\partial \psi}{\partial x}.
\end{equation}

Define-se a variável de similaridade de Blasius:

\begin{equation}
    \eta = y\sqrt{\frac{U}{\nu x}},
\end{equation}

com a função corrente escrita como:

\begin{equation}
    \psi = \sqrt{\nu U x}\,f(\eta).
\end{equation}

A partir dessas definições:

\begin{equation}
    u = U f'(\eta),
\end{equation}

\begin{equation}
    v = \frac{1}{2}\sqrt{\frac{\nu U}{x}}\left[\eta f'(\eta)-f(\eta)\right].
\end{equation}

Substituindo essas expressões na equação de camada limite em $x$, obtém-se a equação de Blasius:

\begin{equation}
    \boxed{f''' + \frac{1}{2}f f'' = 0.}
\end{equation}

As condições de contorno tornam-se:

\begin{align}
    f(0) &= 0, \\
    f'(0) &= 0, \\
    f'(\infty) &= 1.
\end{align}

\section{Sistema equivalente de EDOs de primeira ordem}

Define-se:

\begin{align}
    y_1 &= f, \\
    y_2 &= f', \\
    y_3 &= f''.
\end{align}

Logo:

\begin{align}
    y_1' &= y_2, \\
    y_2' &= y_3, \\
    y_3' &= -\frac{1}{2}y_1y_3.
\end{align}

As condições iniciais são:

\begin{align}
    y_1(0) &= 0, \\
    y_2(0) &= 0, \\
    y_3(0) &= \alpha,
\end{align}

em que $\alpha=f''(0)$ é desconhecido. O valor de $\alpha$ é determinado pelo método de shooting para satisfazer $y_2(\eta_{max})\approx 1$.

\section{Solução numérica por Runge--Kutta}

A integração foi feita pelo método de Runge--Kutta clássico de quarta ordem. Como a condição $f'(\infty)=1$ está no infinito, adotou-se $\eta_{max}=10$, valor suficiente para representar o comportamento assintótico da solução. O parâmetro $\alpha=f''(0)$ foi ajustado pelo método da secante.

O valor numérico obtido foi:

\begin{equation}
    \boxed{f''(0) \approx 0{,}3320573372.}
\end{equation}

Esse valor caracteriza o gradiente de velocidade na parede para a solução de Blasius.

\section{Distribuição de velocidade}

A distribuição de velocidade adimensional é dada por:

\begin{equation}
    \frac{u}{U}=f'(\eta).
\end{equation}

A Figura~\ref{fig:perfil_adimensional} mostra o perfil de velocidade na direção normal à placa, em coordenadas adimensionais.

\begin{figure}[H]
    \centering
    \includegraphics[width=0.78\textwidth]{perfil_velocidade_adimensional.png}
    \caption{Perfil adimensional de velocidade para a camada limite de Blasius.}
    \label{fig:perfil_adimensional}
\end{figure}

Também foi gerado um exemplo dimensional usando $x=1$ m, $U=1$ m/s e $\nu=10^{-6}$ m$^2$/s. Nesse caso:

\begin{equation}
    y = \eta\sqrt{\frac{\nu x}{U}}.
\end{equation}

A estimativa usual para a espessura da camada limite é:

\begin{equation}
    \delta_{99} \approx 5\sqrt{\frac{\nu x}{U}}.
\end{equation}

Para os valores adotados no exemplo:

\begin{equation}
    \delta_{99} \approx 5{,}0\times 10^{-3}\,\text{m}.
\end{equation}

\begin{figure}[H]
    \centering
    \includegraphics[width=0.78\textwidth]{perfil_velocidade_dimensional.png}
    \caption{Exemplo de perfil dimensional de velocidade para $x=1$ m, $U=1$ m/s e $\nu=10^{-6}$ m$^2$/s.}
    \label{fig:perfil_dimensional}
\end{figure}

\section{Código em Python}

O programa em Python resolve a equação de Blasius pelo método de shooting com Runge--Kutta de quarta ordem. O código principal está no arquivo \texttt{camada\_limite\_blasius.py}. O trecho essencial do sistema resolvido é:

\begin{lstlisting}[style=pythonstyle]
def rhs(_eta, state):
    y1, y2, y3 = state
    return (y2, y3, -0.5 * y1 * y3)
\end{lstlisting}

O método de shooting ajusta o valor inicial desconhecido $f''(0)$ até que $f'(\eta_{max})=1$.

\section{Conclusão}

A partir das equações de Navier--Stokes, a análise de ordem de grandeza mostrou que, para número de Reynolds elevado, os termos viscosos longitudinais são desprezíveis em comparação com os termos viscosos normais à placa. Com isso, obteve-se o modelo simplificado de camada limite de Prandtl.

A transformação de similaridade de Blasius reduziu o problema de EDPs para uma EDO não linear de terceira ordem. Essa EDO foi transformada em um sistema de três EDOs de primeira ordem e resolvida numericamente por Runge--Kutta de quarta ordem. O resultado obtido, $f''(0)\approx 0{,}3320573372$, é compatível com a solução clássica de Blasius. O perfil de velocidade mostra a transição suave entre a condição de não deslizamento na parede e a velocidade uniforme do escoamento externo.

\end{document}
